\midsectiontitle{Atributos do Universo}

Os atributos definem as capacidades essenciais de um personagem, moldando sua força física, agilidade, resistência, percepção, intelecto e conexão com o divino.  
Cada atributo influencia diretamente testes, combates e interações narrativas.

\noindent


% =========================
% DIV
% =========================
\originlink{att:divinity}
\noindent
\hypertarget{target:divinity}{\textbf{\baskerfamily\large DIV — Divinity (Divindade)}}\\
A ligação do jogador com o Pináculo. DIV mede a proximidade do personagem com entidades superiores, deuses e 
forças cósmicas para permitir canalizar suas forças como bem entender. DIV é o único atributo que tem 
limite máximo sendo ele 100 pontos.


% =========================
% MIG
% =========================

\originlink{att:might}
\noindent
\textbf{\baskerfamily\large MIG — Might (Poder)}\\
A força bruta do corpo e da vontade física. MIG governa o dano em combate corpo a corpo baseado em força bruta, 
a capacidade de empunhar armas pesadas e realizar feitos de pura potência, como arrombar portões, erguer estruturas ou 
subjugar inimigos pela força.


% =========================
% GRA
% =========================
\originlink{att:grace}
\noindent
\textbf{\baskerfamily\large GRA — Grace (Graça)}\\
Expressa velocidade, precisão e harmonia de movimento. GRA afeta ataques precisos, esquivas, acrobacias, 
qualquer ação que exija controle corporal refinado. Personagens com alta Graça parecem mover-se como se guiados por mãos divinas.

% =========================
% FOR
% =========================
\originlink{att:fortitude}
\noindent
\textbf{\baskerfamily\large FOR — Fortitude (Fortitude)}\\
A resistência do corpo a todos maléficios externos. FOR define pontos de vida, resistência a doenças, 
venenos e efeitos prolongados. Um personagem com alta Fortitude sobrevive onde outros caem.

% =========================
% WIS
% =========================
\originlink{att:wisdom}
\noindent
\textbf{\baskerfamily\large WIS — Wisdom (Sabedoria)}\\
Representa conhecimento do mundo, entendimento de energias arcanas e o poder de ataque com as mesmas. 
WIS controla o aprendizado, o uso de magias, rituais e a compreensão de fenômenos ocultos. 
Baixa Sabedoria não implica ignorância total, mas limita profundidade e alcance intelectual para assuntos específicos.

% =========================
% SEN
% =========================
\originlink{att:sense}
\noindent
\textbf{\baskerfamily\large SEN — Sense (Sentido)}\\
Representa a percepção aguçada do mundo. SEN engloba testes de observação, intuição e detecção. 
Personagens atentos percebem o que outros ignoram, sons, mentiras, presenças invisíveis.


\titlesection{E Afinal, o que é Divindade em Overgrown?}

Divindade representa o nível de progressão do personagem dentro do sistema de Overgrown. O valor de DIV é sempre igual ao nível do personagem: um personagem de nível 0 possui DIV 0, um personagem de nível 1 possui DIV 1 e assim sucessivamente, até o limite de DIV 100.

Na criação, em \textbf{DIV 0}, o personagem recebe \textbf{5 Pontos de Talento}. A cada novo nível de DIV, recebe mais \textbf{5 Pontos de Talento}. Esse é o único orçamento concedido pela progressão de nível. Portanto:

\[\text{Pontos de Talento totais}=5\times(\text{DIV}+1)\]

Quando a árvore de talentos estiver implementada, esses mesmos pontos poderão ser gastos em qualquer árvore disponível. As árvores conterão tanto \textbf{Nós de Atributo}, que aumentam MIG, GRA, FOR, WIS ou SEN, quanto \textbf{Nós de Habilidade}, que concedem habilidades. Não existe uma reserva separada para atributos e outra para habilidades.

\begin{rpgboxtips}{Modelo Atual sem Árvore}
Enquanto a árvore ainda não estiver implementada, os 5 Pontos de Talento recebidos por DIV são aplicados diretamente em atributos. As habilidades continuam sendo escolhidas individualmente conforme as regras atuais. Quando a árvore for implementada, esse procedimento será substituído pela compra de Nós de Atributo e Nós de Habilidade, sem conceder pontos adicionais retroativos.
\end{rpgboxtips}

Os pontos podem ser guardados, mas um ponto aplicado não pode ser transferido sem um efeito de redistribuição. DIV não recebe esses pontos: ele aumenta junto do nível. Bônus temporários e equipamentos não consomem Pontos de Talento e não cumprem pré-requisitos de valor bruto permanente.
Ao atingir 100 pontos de Divindade, o personagem alcança o limite máximo de conhecimento que uma mente humana 
é capaz de suportar do Pináculo, tornando-se o mais próximo possível de uma entidade divina.

Uma Divindade elevada não impõe, obrigatoriamente, mudanças no comportamento do personagem. 
Contudo, para fins de interpretação (\textit{roleplay}), seus efeitos narrativos podem ser definidos 
de campanha à campanha.

No universo de Overgrown, é possível adquirir divindade de várias formas, 
seja entrando em contato com vestígios restantes de Energia Primordial, estudando conhecimentos divinos ou 
por meio de Combate. Todas as criaturas e seres tem vestígios de divindades, 
essa energia não se dissipa, e matar algo que contem divindade resulta na distribuição i
gualitária dela para quem participar do abate. A tabela abaixo representa o progresso necessário para avançar cada nível de DIV; seus valores não são Pontos de Talento disponíveis para compra.

\begin{rpgboxtips}{Exemplos de Progressão}
Em DIV 0, o personagem possui 5 Pontos de Talento totais. Em DIV 1, possui 10; em DIV 10, 55; em DIV 20, 105; e em DIV 100, 505. Esses valores são o orçamento total recebido pela progressão, antes de bônus externos.
\end{rpgboxtips}

\begin{center}
\renewcommand{\arraystretch}{1.3}
\begin{tabular}{@{} cc @{}}
\toprule
\rowcolor{gray!15} \textbf{Intervalo de Divindade} & \textbf{Progresso por Nível} \\ 
\midrule
01 -- 10 & 10 pts \\
11 -- 20 & 20 pts \\
21 -- 30 & 40 pts \\
31 -- 40 & 80 pts \\
41 -- 50 & 160 pts \\
\midrule
51 -- 60 & 320 pts \\
61 -- 70 & 640 pts \\
71 -- 80 & 1.280 pts \\
81 -- 90 & 2.560 pts \\
91 -- 100 & 5.120 pts \\
\bottomrule
\end{tabular}
\end{center}



\titlesection{Estatísticas dos jogadores}

\subsection*{Modificador de Atributo}

O valor bruto de um atributo determina sua reserva de dados. O \textbf{Modificador de Atributo (MOD)} representa a parcela numérica usada em dano, recuperação e efeitos que pedem explicitamente um modificador. O MOD não é somado a um teste, salvo quando a própria regra mandar fazê-lo.

Para um atributo maior que 0:

\[\text{MOD de Atributo}=\left\lceil\frac{\text{Atributo}}{5}\right\rceil\]

Um atributo 0 possui MOD 0. Assim, atributos de 1--5 possuem MOD 1; 6--10, MOD 2; 11--15, MOD 3; 16--20, MOD 4; e assim sucessivamente. Quando uma regra pedir o maior MOD entre dois atributos, compare os valores brutos e use apenas o MOD do maior; em empate, o resultado é o mesmo e nenhum dado adicional é concedido.

\begin{rpgboxtips}{Valor Bruto e MOD}
Pré-requisitos escritos como \textbf{MIG 15}, \textbf{GRA 10} ou semelhantes usam o valor bruto permanente do atributo. Bônus temporários não cumprem pré-requisitos. Uma expressão como \textbf{MOD de MIG} usa a fórmula acima.
\end{rpgboxtips}

Além dos atributos, jogadores tem algum recursos que são usado durante sua aventura: HP, IEP, Resistência e Esquiva.

HP, Health Points ou Pontos de Vida, representam a vida total do jogador, o quanto de dano ele consegue aguentar e 
é diretamente atrelado a sua \link{att:fortitude}{Fortitude}, cada ponto de Fortitude garante ao jogador 5 de vida e o jogador começa com 10 pontos de vida base. 
Um jogador com 10 de Fortitude terá 60 de vida, excluindo bônus e reduções de vida provenientes da árvore de talentos 
e de itens equipados.

IEP, Impossible Energy Points ou Pontos de Energia Impossível \link{res:null}{(Null)}, 
representa a capacidade do jogador de fazer feitos impossíveis, como usar Magias e Maldições. 
O IEP é diretamente ligado ao \link{att:wisdom}{Wisdom} do jogador, para cada ponto de Wisdom 
garante 5 pontos de IEP ao jogador, e o jogador começa com 10 pontos de IEP base. Um jogador com 10 de Wisdom terá 60 de IEP, excluindo bônus 
e reduções de IEP provenientes da árvore de talentos e de itens equipados.

Resistência define a capacidade do jogador de resistir a ataques, sua classe de armadura. 
É diretamente baseada também na \link{att:fortitude}{Fortitude} 
do jogador. A cada 8 Pontos de Fortitude o jogador ganha 1 ponto de Resistência, tendo 5 pontos base. 
Um jogador com 24 pontos de Fortitude terá 8 de resistência (5 base + 3 vindo dos Atributos), 
excluindo bônus e reduções provenientes da da árvore de talentos e de itens equipados.

\begin{rpgboxtips}{Resistência em Combate}
Resistência é um atributo passivo: ela determina se o ataque ou magia acerta o jogador, mas não reduz o dano recebido. Resistência, Valor de Bloqueio e Redução de Dano são estatísticas diferentes. Um \textbf{Bloqueio} não aumenta a Resistência: ele reduz o dano de um ataque que já acertou, conforme as regras de Armaduras e Proteções.
\end{rpgboxtips}

Esquiva define a proficiência em desviar do jogador. Ela é baseada diretamente ao atributo de \link{att:grace}{Grace}. 
A cada 7 Pontos de Grace, o jogador ganha 1 ponto de Esquiva. Todos os jogadores tem Esquiva base em 5. 
Um jogador com 21 de Grace, tem 8 de esquiva (5 base + 3 dos Pontos de Atributo), 
excluindo bônus e reduções provenientes da árvore de talentos e de itens equipados. 
Caso o rolamento de ataque seja menor que a esquiva do jogador, todo o dano é negado.

\begin{rpgboxtips}{Esquiva em Combate}
Esquivar de um ataque é considerado uma reação, alvos que esquivam perdem a reação na rodada. 
Falhar na esquiva ainda gasta sua reação.
\end{rpgboxtips}

Pontos de Combate é um recurso que representa a capacidade do jogador de realizar ações em combate, como ataques extras ou usar habilidades de combate.
Um jogador com 10 de Grace e 10 de Might, tem 22 pontos de combate (2 base + 10 de Grace + 10 de Might), excluindo bônus e reduções provenientes da árvore de talentos e de itens equipados.


